\section{Method}
\label{sec:method}

\subsection{Overview}
\label{sec:method:overview}

CSM-SAM keeps the SAM2~\cite{ravi2024sam2} image encoder (ViT-H, frozen) and
most of the mask decoder frozen and introduces three trainable components
totalling approximately $2.7$M parameters.
Figure~\ref{fig:architecture} gives an overview.

\paragraph{Notation.}
Let $\mathbf{F}^{\text{pre}} \in \mathbb{R}^{N \times h \times w \times C}$
denote the SAM2 encoder features for the $N$ pre-RT slices,
$\mathbf{F}^{\text{mid}} \in \mathbb{R}^{h \times w \times C}$ the
current mid-RT slice features, $t \in \mathbb{R}_{>0}$ the elapsed weeks,
and $y \in \{0,1\}^{H \times W}$ the ground-truth mid-RT mask.

\subsection{Cross-Session Memory Attention}
\label{sec:method:csma}

A \texttt{CrossSessionMemoryEncoder} projects the pre-RT features into a
compact memory bank $\mathbf{M}_{\text{pre}} \in \mathbb{R}^{N_{\text{mem}} \times C}$
via spatial pooling and mask-guided modulation.  A
\texttt{CrossSessionMemoryAttention} module then lets the mid-RT query tokens
attend to $\mathbf{M}_{\text{pre}}$ through standard multi-head cross-attention:
\begin{equation}
  \mathbf{F}^{\text{cross}} = \mathrm{Attn}\!\left(
    \mathbf{F}^{\text{mid}},\;
    \mathbf{M}_{\text{pre}} + \phi(t),\;
    \mathbf{M}_{\text{pre}}
  \right),
\end{equation}
where $\phi(t)$ is a sinusoidal+MLP continuous-time embedding of weeks elapsed.
A learned sigmoid gate blends $\mathbf{F}^{\text{cross}}$ with within-session
memory (prior mid-RT slices) to combine cross-visit and intra-volume context.

\subsection{Variational Change Latent}
\label{sec:method:latent}

The core novel contribution is a variational module that infers a
low-dimensional latent code $\mathbf{z} \in \mathbb{R}^{d_z}$ representing the
patient's \emph{treatment response phenotype}: a compact summary of how this
patient's tumour changes under radiotherapy.

\paragraph{Training (posterior inference).}
At training time, both sessions are available.  A
\texttt{ChangeLatentEncoder} computes the approximate posterior:
\begin{equation}
  q(\mathbf{z} \mid \mathbf{f}_{\text{pre}},\, \mathbf{f}_{\text{mid}},\, t)
  = \mathcal{N}\!\left(\boldsymbol{\mu}_q,\, \mathrm{diag}(\boldsymbol{\sigma}_q^2)\right),
\end{equation}
where $\mathbf{f}_{\text{pre}} = \mathrm{GAP}(\mathbf{M}_{\text{pre}})$ and
$\mathbf{f}_{\text{mid}} = \mathrm{GAP}(\mathbf{F}^{\text{mid}})$ are
global-average-pooled summaries.
The posterior mean and log-variance are computed by a three-layer MLP:
$[\mathbf{f}_{\text{pre}},\,\mathbf{f}_{\text{mid}},\,\phi_z(t)] \to
(\boldsymbol{\mu}_q,\, \log\boldsymbol{\sigma}_q^2)$,
where $\phi_z$ is a separate sinusoidal temporal encoder.
A reparameterised sample $\mathbf{z} = \boldsymbol{\mu}_q + \boldsymbol{\sigma}_q \odot \boldsymbol{\varepsilon}$,
$\boldsymbol{\varepsilon} \sim \mathcal{N}(\mathbf{0}, \mathbf{I})$, is used
for all downstream computations.

\paragraph{Inference (prior + retrieval).}
At test time, the mid-RT scan has not yet been acquired.  A
\texttt{ChangePrior} network estimates the prior distribution:
\begin{equation}
  p(\mathbf{z} \mid \mathbf{f}_{\text{pre}},\, t)
  = \mathcal{N}\!\left(\boldsymbol{\mu}_p,\, \mathrm{diag}(\boldsymbol{\sigma}_p^2)\right).
\end{equation}
The prior is further refined by cross-patient retrieval: the top-$K$
training patients most similar to the query (cosine similarity on
$\mathbf{f}_{\text{pre}}$) contribute their stored posterior means
$\{\mathbf{z}_i\}_{i=1}^K$ as an empirical estimate of the posterior:
\begin{equation}
  \mathbf{z}_{\text{ret}} = \sum_{i=1}^{K} w_i\, \mathbf{z}_i,
  \qquad
  w_i \propto \exp\!\left(\frac{\mathrm{sim}(\mathbf{f}_{\text{pre}},\, \mathbf{f}_{\text{pre}}^{(i)})}{\tau}\right),
\end{equation}
and the final test-time latent is the blend
$\mathbf{z} = (1-\alpha)\,\boldsymbol{\mu}_p + \alpha\,\mathbf{z}_{\text{ret}}$
with $\alpha = 0.5$, $\tau = 0.1$.

\paragraph{FiLM conditioning.}
The latent $\mathbf{z}$ conditions the cross-session attention query via
Feature-wise Linear Modulation~\cite{perez2018film}:
\begin{equation}
  \mathbf{q}' = \mathbf{q} \odot \bigl(1 + \boldsymbol{\gamma}(\mathbf{z})\bigr)
              + \boldsymbol{\beta}(\mathbf{z}),
\end{equation}
where $\boldsymbol{\gamma}, \boldsymbol{\beta} \in \mathbb{R}^C$ are produced
by a two-layer MLP.  This makes the attention pattern \emph{change-type-aware}:
a fast-responder $\mathbf{z}$ retrieves different memory tokens from
$\mathbf{M}_{\text{pre}}$ than a non-responder $\mathbf{z}$.  Both heads are
zero-initialised, so the contribution is near-zero at the start of training
and grows as the latent space is shaped.

\subsection{Auxiliary Heads}
\label{sec:method:auxiliary}

A \texttt{ChangeHead} predicts a four-class change map
(background / stable / grown / shrunk) supervised by the signed XOR of the
pre- and mid-RT masks.  A dual-head decoder invokes the SAM2 mask decoder
once per structure (GTVp, GTVn), conditioned on the corresponding pre-RT mask
as a dense prompt.

\subsection{Training Objective}
\label{sec:method:loss}

\begin{equation}
  \mathcal{L} =
    \mathcal{L}_{\text{Dice}} + \mathcal{L}_{\text{BCE}}
    + \lambda_c\, \mathcal{L}_{\text{change}}
    + \beta(e)\cdot\lambda_{\text{KL}}\, \mathcal{L}_{\text{KL}},
\end{equation}
where $\lambda_c = 0.3$, $\lambda_{\text{KL}} = 0.1$, and
\begin{equation}
  \mathcal{L}_{\text{KL}} = \mathrm{KL}\!\left(
    q(\mathbf{z} \mid \mathbf{f}_{\text{pre}},\, \mathbf{f}_{\text{mid}},\, t)
    \;\|\;
    p(\mathbf{z} \mid \mathbf{f}_{\text{pre}},\, t)
  \right).
\end{equation}
Critically, the KL is taken against the \emph{learned prior} $p$, not a
standard normal.  This forces $p$ to model the population-level distribution
of treatment responses, while $q$ is regularised toward it.  The annealing
schedule $\beta(e) = \min(1,\, e / E_{\text{warm}})$ with
$E_{\text{warm}} = 10$ epochs prevents the KL from collapsing the posterior
before the segmentation loss has shaped the latent space.

\subsection{Emergent Phenotype Discovery}
\label{sec:method:phenotype}

No phenotype labels appear anywhere in the training signal.  The only
supervision is the segmentation objective and the KL regulariser.  The
claim tested in Section~\ref{sec:results:phenotype} is that $\boldsymbol{\mu}_q$
nonetheless organises into clusters that correlate with the volume-change ratio
$\Delta V = (\mathrm{vol}(m_{\text{mid}}) - \mathrm{vol}(m_{\text{pre}})) /
\mathrm{vol}(m_{\text{pre}})$, a clinically meaningful indicator of treatment
response.  This is verified post-hoc via Spearman rank correlation on held-out
patients, not through any label in the training set.
